% ===============================
% Worksheet / Manual Template
% ===============================
\documentclass[12pt, a4paper]{article}

% ===============================
% Metadata
% ===============================
\newcommand*{\CourseCode}{MAT321}
\newcommand*{\CourseTitle}{Real Analysis II}
\newcommand*{\Chapter}{Sequential Compactness}
\newcommand*{\Topics}{Finite Intersection Property, Bolzano-Weierstrass Property}
\newcommand*{\DocDate}{April 25, 2026}

\include{header}

% ==========================================
% Document
% ==========================================
\begin{document}
\FrontPage
\Large

\begin{heading}
    Finite Intersection Property (FIP)
\end{heading}

\vspace{10pt}

\begin{definition}[Finite Intersection Property]
    A family of sets $\{A_\alpha\}_{\alpha \in I}$ is said to have the \textit{finite intersection property} if for every finite subset $J \subseteq I$, we have
    \[
        \bigcap_{\alpha \in J} A_\alpha \neq \emptyset.
    \]
\end{definition}

\vspace{20pt}

\begin{problem}[Example of FIP]
    Consider the family of sets $\{A_n\}_{n=1}^\infty$ where $A_n = \left[-\frac{1}{n}, \frac{1}{n}\right]$. Show that this family has the finite intersection property, but the intersection of all sets in the family is empty.
    
\end{problem}

\clearpage

\begin{problem}[Example of FIP]
    Consider the family of sets $\{A_n\}_{n=1}^\infty$ where $A_n = \left[n, \infty\right)$. Show that this family has the finite intersection property, but the intersection of all sets in the family is empty.
    
\end{problem}

\clearpage

\begin{heading}
    Compactness and Finite-Intersection Property
\end{heading}

\vspace{10pt}

\begin{theorem}[Compactness and FIP]
    A subset $K$ of a metric space is compact if and only if for every family of closed subsets $\{F_\alpha\}_{\alpha \in I}$ of $K$ with the finite intersection property, we have
    \[
        \bigcap_{\alpha \in I} F_\alpha \neq \emptyset.
    \]
\end{theorem}

\clearpage

\begin{heading}
    Sequential Compactness and Bolzano-Weierstrass Property
\end{heading}

\vspace{10pt}

\begin{definition}[Sequential Compactness]
    A subset $K$ of a metric space is said to be \textit{sequentially compact} if every sequence in $K$ has a convergent subsequence whose limit is in $K$.
\end{definition}

\clearpage

\begin{definition}[Bolzano-Weierstrass Property]
    A subset $K$ of a metric space is sequentially compact if and only if every infinite subset of $K$ has a limit point in $K$.
\end{definition}

\clearpage

\begin{theorem}[Sequential Compactness and B-W Property]
    A subset $K$ of a metric space is sequentially compact if and only if it has the Bolzano-Weierstrass property.
\end{theorem}

\clearpage

\begin{heading}
    Compactness and Sequential Compactness
\end{heading}

\vspace{10pt}

\begin{theorem}[Compactness implies Sequential Compactness]
    Every compact metric space is sequentially compact.
\end{theorem}

\clearpage

\begin{theorem}[Sequential Compactness implies Compactness]
    Every sequentially compact metric space is compact.
\end{theorem}

\clearpage

\begin{problem}[Remark on Compactness and Sequential Compactness]
    In general, compactness and sequential compactness are not equivalent in arbitrary topological spaces. However, in metric spaces, they coincide. Provide an example of a topological space where compactness does not imply sequential compactness.
\end{problem}




\EndPage
\end{document}